\documentclass[instructions.tex]{subfiles}
\begin{document}
 
\chapter{De la fuite du monde.}
 
Jésus-Christ nous ordonne de fuir le monde, parce que l'esprit du monde, les conversations du monde, les divertissemens du monde sont les trois grands écueils du salut.
 
\primo Le Sauveur, en nous ordonnant de fuir le monde, ne nous défend pas d'habiter dans le monde; mais il nous défend de vivre selon l'esprit du monde, et d'aimer les choses et les vanités du monde, parce que celui qui s'y attache et qui les aime, se rend ennemi de Dieu\footnote{Amicitia hujus mundi inimica est Dei. Jac. 4.}. Le monde est tout plongé dans la malignité, dit saint Jean. Son esprit est incompatible avec l'esprit de Dieu\footnote{Quem mundus non potest accipere. Joan. 14.}. C'est un esprit de volupté, d'avarice et d'intérêt, un esprit de jalousie, de fourberie, de duplicité et d'orgueil, qui n'inspire que l'amour du plaisir et des richesses, l'ambition et la discorde.
 
En fréquentant le monde, on prend son esprit, on goûte ses maximes, on voit ses exemples, et on les suit. De là tant de désordres et de crimes en tout genre; de là le dégoût des choses de Dieu, les inquiétudes qui, au milieu des embarras du siècle, et même au milieu des folles joies du monde, tourmentent la conscience.
 
C'est donc chez vous, plutôt qu'ailleurs, et dans les occupations de votre état que vous devez vous plaire. Une femme qui est dans son ménage; un père, dans sa famille et dans ses emplois; un homme d'affaires, dans son bureau; un homme d'étude, dans son cabinet; l'enfant, le domestique, sous les yeux de ses parens, de ses maîtres; l'ouvrier, dans son travail, sont à couvert des grands pièges du monde. Partout ailleurs on est hors de son élément, comme le poisson hors de l'eau; on y respire un air meurtrier.
 
Qui veut se sauver, trouve les compagnies ennuyeuses. Il trouve sa consolation dans la prière et dans le lieu saint, et sa sûreté, dans une vie retirée. Moins il y aura en vous de curiosité et d'épanchement au-dehors, plus Dieu se fera sentir au-dedans.
 
Rien ne doit vous attirer dans les compagnies du monde, que la nécessité, la bienséance, l'obéissance ou la charité. Quiconque y paraît sans quelques-unes de ces raisons, se met dans le danger; mais souvenez-vous, quand vous êtes seul, de ne point faire ce que vous voudriez cacher aux hommes, et de ne point penser à tout ce qui déplaît à Dieu.
 
\secundo Quant à la conversation, elle est inévitable dans le monde; mais c'est une science aussi admirable qu'elle est rare, de converser avec les hommes sans offenser Dieu, et sans se nuire à soi-même. On converse dès l'enfance, et à l'âge de soixante ans à peine sait-on converser utilement. C'est avec raison qu'un auteur a dit : Autant de fois que je me suis rencontré parmi les hommes, j'en suis revenu moins homme, c'est-à-dire moins raisonnable, et moins uni à Dieu.
 
Ne vous attachez pas à quantité d'amis, et ne vous flattez pas d'en avoir beaucoup, parce qu'il y en a peu de véritables. Choisissez-les, dit le Sage, entre mille, qui soient gens de bien, chastes et sincères; de tels amis sont un trésor. Pour ne pas vous y tromper, fréquentez les plus vertueux, afin que vous ne voyiez dans vos amis que ce qui peut vous rendre meilleur.
 
Pour vivre agréablement et saintement avec vos amis et avec tout le monde, ayez toujours cette humble complaisance qui sait s'accommoder à l'humeur des autres. Entrez dans leur sentiment, pourvu que la conscience n'y soit pas intéressée. Dans la conversation, ne faites peine à personne, défendez les absens, épargnez vos ennemis, ménagez vos amis. S'il est défendu de faire peine à un ennemi, à plus forte raison, à un ami. Railler les autres, c'est se rendre méprisable; parler des services qu'on a rendus, c'est les faire acheter; en faire des reproches, c'est manquer d'éducation.
 
Quoiqu'on doive être réservé en conversation, il faut néanmoins y être affable, et ne gêner personne; mais il est dangereux d'être enjoué à l'égard des personnes du sexe. C'est dans ces occasions qu'il faut, comme dit le Sage, entourer ses oreilles d'épines, et veiller sur ses regards. Ce qu'on entend et ce qu'on voit pénètre souvent jusqu'au fond de l'âme, et la souille.
 
Le poison se glisse dans le cœur, surtout dans les promenades et dans les veillées avec le sexe, par les conversations trop libres, et par les entretiens peu chastes. Plus ce poison paraît agréable, plus il est mortel. C'est en ce point, plus qu'en tout autre, que celui qui méprise les moindres choses, tombera peu à peu dans de grands désordres\footnote{Qui spernit modica, paulatim decidet. Eccl.}. Tel se croit innocent, qui, devant Dieu, est déjà souillé.
 
\section{Résolutions}
 
\primo Puisque le monde offre de si nombreux dangers, et qu'il est très-difficile d'y conserver l'esprit de Jésus-Christ, je l'éviterai autant que les devoirs de mon état, les bienséances et la charité me le permettront.

\secundo Quand je serai forcé de m'y trouver, je m'y rendrai après avoir invoqué le secours de la grâce, et m'être mis en garde contre la séduction, et j'y conserverai plus que je pourrai le souvenir de la présence de Dieu, pensant souvent qu'il me voit.

\tertio Enfin, je ne mettrai au nombre de mes amis, et je ne fréquenterai que des personnes avec lesquelles je ne courrai aucun danger de me perdre.

\prayer{Mon Dieu, donnez-moi de l'éloignement pour toutes les compagnies où je me verrais exposé à vous offenser.}
 
\end{document}
